\chapter*{Resumen}
\addcontentsline{toc}{chapter}{Resumen}

WebNLG es uno de los \textit{benchmarks} de referencia para la generación de
lenguaje natural a partir de tripletas RDF de DBpedia. A lo largo de sus
sucesivas ediciones (2017, 2020 y 2023) el corpus se ha publicado en inglés,
ruso y un conjunto de lenguas con recursos limitados, pero no en español. Este
Trabajo de Fin de Máster aborda esa ausencia con el objetivo de crear la
infraestructura de soporte necesaria para producir y verificar una eventual
extensión al español del corpus.

La aportación se articula en torno a tres ejes. En primer lugar, se diseña y se
desarrolla \texttt{lanbench}, una plataforma web propia que materializa el
flujo \textit{Human-in-the-Loop} de anotación y revisión. La aplicación se
construye como un monolito modular por capas sobre Node.js, Express, Prisma y
MariaDB, despliega sus componentes mediante Docker Compose y articula cuatro
roles funcionales --anotador, revisor, administrador y agente de IA-- sobre un
modelo de \textit{datasets}, permisos y estadísticas que controla la
concurrencia y la trazabilidad. En segundo lugar, se integra un componente de
modelos de lenguaje de gran escala (LLM) en dos modalidades operativas: un
\textit{corrector} que detecta y clasifica errores ortográficos, gramaticales y
semánticos sobre frases producidas por el anotador, y un \textit{creativo} que
genera verbalizaciones a partir de las tripletas RDF, dejando al humano el
papel de editor. En tercer lugar, se diseña una experimentación reproducible
que cuantifica la calidad de ambos flujos sobre varios conjuntos de prueba, con
dos proveedores comerciales (Groq Llama 3.3 70B y Google AI Studio Gemini 2.5
Flash) y bajo configuración idéntica del \textit{prompt}.

Los resultados muestran que con un ajuste dirigido de léxico determinista en
\texttt{domain/\allowbreak spanish/\allowbreak lexicon.js} la generación de Llama 3.3 70B mejora el
acuerdo con la referencia humana del 60\,\% al 85\,\% sobre el corpus de 20
entradas. Además, sobre el corpus extendido de 50 entradas se obtiene que Llama
3.3 70B obtiene un acuerdo del 68\,\%, con desacuerdos concentrados en la
falta de fluidez. En cambio, la submuestra efectiva de Gemini alcanza el 91\,\%
de acuerdo, sugiriendo una mayor competencia que Llama 3.3 70B en este
\textit{pipeline}. El flujo de generación queda validado de extremo a extremo,
si bien la cuantificación a escala real queda enormemente condicionada por las
cuotas gratuitas de los proveedores.

\vspace{1em}
\noindent\textbf{Palabras clave:} WebNLG, RDF, generación de lenguaje natural,
\textit{Human-in-the-Loop}, modelos de lenguaje de gran escala,
\textit{crowdsourcing} académico, métricas de evaluación, español.

%%--------------
\newpage
%%--------------

\chapter*{Abstract}
\addcontentsline{toc}{chapter}{Abstract}

WebNLG is one of the reference \textit{benchmarks} for natural language
generation from DBpedia RDF triples. Across its successive editions (2017, 2020
and 2023) the corpus has been published in English, Russian and a set of
low-resource languages, but not in Spanish. This Master's Thesis addresses this
gap with the aim of building the supporting infrastructure needed to produce
and verify a potential Spanish extension of the corpus.

The contribution is organised around three axes. First, \texttt{lanbench} is
designed and developed, a custom web platform that materialises a
\textit{Human-in-the-Loop} annotation and review workflow. The application is
built as a modular, layered monolith on top of Node.js, Express, Prisma and
MariaDB, deploys its components through Docker Compose and articulates four
functional roles --annotator, reviewer, administrator and AI agent-- on top of
a dataset, permission and statistics model that controls concurrency and
traceability. Second, a Large Language Model (LLM) component is integrated in
two operational modes: a \textit{corrector} that detects and classifies
orthographic, grammatical and semantic errors on sentences produced by the
annotator, and a \textit{creative} mode that generates verbalisations from the
RDF triples, leaving the human in the role of editor. Third, a reproducible
experiment is designed that quantifies the quality of both flows across several
test sets, with two commercial providers (Groq Llama 3.3 70B and Google AI
Studio Gemini 2.5 Flash) and under identical \textit{prompt} configuration.

The results show that, with a targeted tuning of the deterministic lexicon at
\texttt{domain/\allowbreak spanish/\allowbreak lexicon.js}, the agreement of Llama 3.3 70B generation
with the human reference improves from 60\,\% to 85\,\% on the 20-entry corpus.
Furthermore, on the extended 50-entry corpus Llama 3.3 70B obtains a 68\,\%
agreement, with disagreements concentrated on the lack of fluency. By contrast,
the effective Gemini sub-sample reaches a 91\,\% agreement, suggesting a
greater competence than Llama 3.3 70B on this \textit{pipeline}. The generation
flow is validated end-to-end, although real-scale quantification remains
heavily constrained by the providers' free-tier quotas.

\vspace{1em}
\noindent\textbf{Keywords:} WebNLG, RDF, Natural Language Generation,
\textit{Human-in-the-Loop}, Large Language Models, academic
\textit{crowdsourcing}, evaluation metrics, Spanish.

%%%%%%%%%%%%%%%%%%%%%%%%%%%%%%%%%%%%%%%%%%%%%%%%%%%%%%%%%%%
%% Final del resumen.
%%%%%%%%%%%%%%%%%%%%%%%%%%%%%%%%%%%%%%%%%%%%%%%%%%%%%%%%%%%
